% Generated draft: requires mathematical and source review.
\documentclass[11pt]{article}
\usepackage[margin=25mm]{geometry}
\usepackage{fontspec}
\setmainfont{Latin Modern Roman}
\usepackage{amsmath,amssymb,mathtools,unicode-math}
\setmathfont{Latin Modern Math}
\usepackage{enumitem,needspace,xurl,hyperref}
\setlength{\parindent}{0pt}
\setlength{\parskip}{5pt}
\setlength{\emergencystretch}{3em}
\newcommand{\N}{\mathbb{N}}
\newcommand{\Z}{\mathbb{Z}}
\newcommand{\Q}{\mathbb{Q}}
\newcommand{\R}{\mathbb{R}}
\newcommand{\C}{\mathbb{C}}
\newcommand{\F}{\mathbb{F}}
\newcommand{\A}{\mathbb{A}}
\newcommand{\PP}{\mathbb{P}}
\newcommand{\E}{\mathbb{E}}
\newcommand{\eps}{\varepsilon}
\newcommand{\dd}{\,\mathrm{d}}
\begin{document}
\section*{UnsolvedMath ID: 701185}
\textbf{Source problem code:} catalogue entry
\subsection{Definitions and mathematical statement}
Fix a prime \(p\). A group \(G\) is locally finite if every finitely generated subgroup of \(G\) is finite. A \(p\)-element is an element of order \(p^k\) for some integer \(k\geq 0\). A \(p'\)-group here means a periodic group in which every element has order relatively prime to \(p\).

Following the definition attributed in the source to S. N. Chernikov, \(G\) satisfies the \(p\)-minimal condition if there is no infinite strictly descending chain
\[
H_1>H_2>H_3>\cdots
\]
of subgroups of \(G\) such that, for every \(i\geq 1\), the set \(H_i\setminus H_{i+1}\) contains a \(p\)-element.

For a prime \(q\), the quasicyclic group \(C_{q^\infty}\) is the union of cyclic groups of orders \(q,q^2,q^3,\ldots\) under the usual inclusions. A group is Chernikov if it has a subgroup of finite index isomorphic to a finite direct product
\[
\prod_{j=1}^{m} C_{q_j^\infty},
\qquad m\geq 0,
\]
where the \(q_j\) are primes. The empty product is permitted, so finite groups are included.

Under the finite-length convention, a subnormal series is a chain
\[
1=G_0\trianglelefteq G_1\trianglelefteq\cdots\trianglelefteq G_r=G,
\qquad r\geq 0,
\]
where each \(G_{i-1}\) is normal in \(G_i\). The supplied source does not explicitly state whether its term ``subnormal series'' permits transfinite series; the displayed formulation uses the finite-length convention, and any intended broader interpretation remains unresolved.

\textbf{Question.} For every prime \(p\) and every locally finite group \(G\) satisfying the \(p\)-minimal condition and admitting such a subnormal series with each factor \(G_i/G_{i-1}\) finite or a \(p'\)-group, must
\[
P_p(G):=\bigl\langle x\in G:\operatorname{ord}(x)=p^k\text{ for some integer }k\geq 0\bigr\rangle
\]
be a Chernikov subgroup of \(G\)? The conclusion concerns the subgroup generated by all \(p\)-elements; it does not additionally require that this subgroup be a \(p\)-group.
\subsection{Short English statement}
Does the subgroup generated by all \(p\)-elements of a locally finite group have to be Chernikov when the group satisfies the specified descending-chain condition and has a subnormal series with factors that are finite or \(p'\)-groups? The draft explicitly uses a finite subnormal series; whether the source intends a broader series convention requires review.
\subsection{Sources}
\begin{itemize}
\item Local UnsolvedMath registry, supplied complete source record: permanent ID \(701185\), exact source problem code \texttt{catalogue entry}, title \emph{Kourovka Notebook Problem 15.97}. The supplied registry status is \texttt{open}; this label is not an independently verified literature assessment.
\item \emph{The Kourovka Notebook (Unsolved Problems in Group Theory)}, Problem 15.97, attributed to N. S. Chernikov, proposed in 2002. The supplied background identifies the 15th issue (2002) and the 21st-issue edition (2026), \texttt{arXiv:1401.0300v48}. The supplied collection description names E. I. Khukhro and V. D. Mazurov as editors. These bibliographic details are supplied attribution and have not been independently checked.
\item The source credits S. N. Chernikov with the definition of the \(p\)-minimal condition; this attribution is distinct from the problem's attribution to N. S. Chernikov.
\end{itemize}
\end{document}
